\chapter{Conclusions}
\label{conclusion}

The primary objective of this thesis was to design and implement a scheduler capable of dynamically assigning computational tasks, specifically streams of matrix multiplications, to the different accelerators found in modern architectures. The goal was to overcome the limitations of traditional software, which typically relies on a single processor, by simultaneously exploiting the entire heterogeneous topology of a computer: the CPU, the Integrated GPU, the Discrete GPU, and the Neural Processing Unit.

To achieve this, an essential preliminary phase focused on integrating and studying different compute backends (CUDA, SYCL, OpenVINO, and OpenBLAS) to create a uniform abstraction layer over the hardware. Following this, the AMM Scheduler was developed, featuring distinct scheduling logics designed to orchestrate these asymmetric resources in real time.

The experimental results, conducted on a complex hybrid architecture, validated the effectiveness of the proposed solution. The system successfully demonstrated that coordinating multiple accelerators can lead to performance gains and increased energy efficiency.

Specifically, the comparison between the scheduling logics highlighted distinct behaviors depending on the workload:

\begin{itemize}
\item \textbf{Homogeneous Workloads}: For streams of identical size (1024x1024 and 2048x2048), the dynamic strategy proved to be effective. Its greedy nature allowed it to minimize idle time by keeping all accelerators busy, achieving speedups of up to 2.20 compared to the single CUDA card. However, as the matrix size increased to heavy workloads, the static strategy demonstrated superior performance (2x speedup vs 1.78x for dynamic). The predictive nature of the static logic prevented, in most cases, a slow device that delays the entire batch by taking on a final massive task.

\item \textbf{Heterogeneous Workloads}: In scenarios with mixed matrix sizes (streams from 512x512 to 4096x4096), the necessity of a smart, predictive scheduler became evident. The static hetero strategy successfully mapped task complexity to hardware capability, achieving 1.65x speedup. In contrast, the dynamic logic failed to handle the variance, resulting in severe performance degradation (0.42x speedup) due to inefficient task assignment, proving that a simple greedy approach is insufficient for diverse workloads.

\item \textbf{Large Matrix Splitting}: The split strategy demonstrated that even for a single massive operation, decomposing the problem and offloading chunks to the CPU and iGPU can yield a 1.49x speedup and reduce total energy consumption by 21\%.
\end{itemize}

In conclusion, all the initial project objectives were met. The AMM Scheduler proved that an approach to hardware that embraces the diversity of available accelerators rather than ignoring it, is a viable path to maximizing both throughput and energy efficiency in modern computing systems.

\section{Future Work}
While the current implementation offers the initial foundation for heterogeneous scheduling, several avenues for future research and optimization remain open to enhance this project capabilities:

\begin{itemize}
\item Multi Node: The evolution of this project could be to extend the scheduler's scope beyond a single machine. Future work could focus on adapting the coordinator to distribute workloads across a server cluster, treating remote nodes as additional accelerators. This would involve implementing network aware logic to balance computation with network latency, and hardware aware if the nodes have all different accellerators. 

\item Expansion of Supported Primitives and Logics: Currently the scheduler is specialized in Matrix Multiplication. Future developments could aim to extend the library of supported operations, adding other primitives such as Convolutions, FFTs, and so on. Moreover the current logic handles streams of independent tasks. A significant step would be the implementation of a dependency aware scheduler capable of managing Directed Acyclic Graphs (DAGs). This would allow the system to orchestrate entire computational pipelines like a full Neural Network inference.
\end{itemize}

\section{Personal Balance}
This thesis represented my first encounter with a software project of this magnitude and complexity. It has been a long journey, characterized by numerous challenges. One of the most significant aspects was the responsibility of making critical architectural decisions early on, choices that would dictate the entire development of the project. The uncertainty of whether these foundational decisions would prove effective could only be verified upon completion. 
However, this pressure ultimately served as a driving force. It compelled me to approach all the challenges with a pragmatic mindset, and in the end the results validated the effort.

Above all, as a hardware enthusiast, working on this project was very exciting. I had the chance to interact directly with many different accelerators and orchestrate them, exploring and exploiting all of their libraries. This experience taught me a lot about their architectures and gave me a strong foundation in High Performance Computing.

I would like to express my sincere gratitude to my supervisor for the trust placed in me for this thesis. I deeply appreciate the freedom granted to explore this project, accompanied by his constant guidance and ability to point me in the right direction whenever necessary. Confirming what I already knew, and as this experience has fully confirmed, it has been a true pleasure and privilege to work on this Master's thesis with him.

In conclusion, I remain profoundly satisfied with this work. It has fostered my growth both personally and as a software engineer, fueling my ambition to contribute to a better (and faster!) world.

